% Generated draft: requires mathematical and source review.
\documentclass[11pt]{article}
\usepackage[margin=25mm]{geometry}
\usepackage{fontspec}
\setmainfont{Latin Modern Roman}
\usepackage{amsmath,amssymb,mathtools,unicode-math}
\setmathfont{Latin Modern Math}
\usepackage{enumitem,needspace,xurl,hyperref}
\setlength{\parindent}{0pt}
\setlength{\parskip}{5pt}
\setlength{\emergencystretch}{3em}
\newcommand{\N}{\mathbb{N}}
\newcommand{\Z}{\mathbb{Z}}
\newcommand{\Q}{\mathbb{Q}}
\newcommand{\R}{\mathbb{R}}
\newcommand{\C}{\mathbb{C}}
\newcommand{\F}{\mathbb{F}}
\newcommand{\A}{\mathbb{A}}
\newcommand{\PP}{\mathbb{P}}
\newcommand{\E}{\mathbb{E}}
\newcommand{\eps}{\varepsilon}
\newcommand{\dd}{\,\mathrm{d}}
\begin{document}
\section*{UnsolvedMath ID: 30003520}
\textbf{Source problem code:} catalogue entry
\subsection{Definitions and mathematical statement}
\textbf{Source question.} Are the observed spectral features---nontrivial discrete spectrum and, in irrational geometry, existence of a frequency limit point---general for nonlinear boundary-value problems with energy-dependent boundary conditions?

\textbf{Terminology requiring specification.} Once a particular boundary-value problem, admissible solution class, spectral parameter space \(\Lambda\), and topology on \(\Lambda\) have been specified, let \(\Sigma\subseteq\Lambda\) denote its admissible spectral parameters. A point \(\lambda\in\Sigma\) is isolated if some open neighborhood \(U\) of \(\lambda\) satisfies \(U\cap\Sigma=\{\lambda\}\). A point \(\lambda_*\in\Lambda\) is an accumulation point of \(\Sigma\) if every open neighborhood \(U\) of \(\lambda_*\) satisfies
\[
U\cap\bigl(\Sigma\setminus\{\lambda_*\}\bigr)\ne\varnothing.
\]
These are conditional definitions, not additional hypotheses supplied by the record. The source does not specify whether its discrete spectrum means isolated admissible parameters, isolated operator eigenvalues with additional multiplicity conditions, or another nonlinear spectral notion. It also does not identify frequency with the energy parameter.

\textbf{Unresolved formulation.} The record asks whether the two observed features persist in a broader class, with the frequency-limit-point feature qualified by irrational geometry. It does not specify the nonlinear equation, domain or graph, boundary conditions, admissible functions, energy or frequency range, geometric ratio whose irrationality is relevant, or class over which generality is intended. Consequently, it does not determine a precise universally quantified conjecture. In particular, neither the existence of isolated spectral parameters nor the existence of a frequency accumulation point is asserted here as a theorem. The supplied background relates the question to double-bridge standing waves and to an unresolved complete classification of an original energy-dependent boundary-value problem, but does not provide that problem's equations or establish that classification as equivalent to the stated generality question.
\subsection{Short English statement}
The supplied question asks how broadly two spectral phenomena persist for nonlinear boundary-value problems with energy-dependent boundary conditions: a nontrivial discrete spectrum and, when the geometry is irrational, a frequency accumulation point. The equations, admissible class, spectral definitions, and intended quantifiers are missing, so a precise conjecture requires further source clarification. The supplied open label and literature assessment are dated curation, not independently verified conclusions.
\subsection{Sources}
\begin{itemize}
\item Supplied local UnsolvedMath registry record, permanent ID \(30003520\), exact source problem code \texttt{catalogue entry}, title: \emph{Spectra of Nonlinear Problems with Energy-Dependent Boundary Conditions}. The record lists proposed year 2017 and provides no named proposer.
\item Supplied report citation: \emph{Nonlinear Partial Differential Equations on Graphs (2018)}, Oberwolfach Reports, pp.\ 1805--1868, DOI \texttt{10.4171/owr/2017/29}. \url{https://doi.org/10.4171/owr/2017/29}. The background describes Oberwolfach Reports as documenting workshops and their open-problem discussions; the literature-source label is ``Oberwolfach Report 29/2017.'' The supplied title year and report designation are retained without resolving their discrepancy.
\item Supplied literature reference: Noja--Rolando--Secchi, \emph{Journal of Differential Equations}. \url{https://doi.org/10.1016/j.jde.2018.07.038}. The supplied assessment calls this a 2019 follow-up, says it classifies a substantial subfamily of double-bridge standing waves, and says the complete classification of the original energy-dependent boundary-value problem remains unresolved. These claims have not been independently verified.
\item Literature assessment date, copied exactly: \texttt{2026-08-21}. The supplied dated curation labels the problem open and lists the preceding journal article and Oberwolfach report as sources checked. This records the supplied assessment only; it does not document independent consultation or verification of either source for this draft.
\end{itemize}
\end{document}
